\section{Exam Strategy and Mixed Task Recipes}

\subsection{Ctrl+F keywords: multiple choice, statement validity, units, limiting case, answer check, traps, Danish keywords}
Use for cross-topic exam tactics that recur in marked multiple-choice and calculation problems.

\subsection{Multiple-choice statement elimination}

\subsubsection{One false clause makes an option false}
\[
\text{valid answer option}=\text{all required clauses true}
\]
For statement options, test each clause independently: sign, unit, condition, and exception. Eliminate options with one wrong condition even if the formula name looks right.

\subsubsection{Dimensional and unit check}
\[
\text{final unit must match requested quantity}
\]
Use units as algebraic factors. Common exam errors are Celsius in gas/Arrhenius equations, entropy in \(\mathrm J\) with enthalpy in \(\mathrm{kJ}\), and cell length units in density.

\subsubsection{Logarithm checks for pH, \(K\), and Nernst}
\[
\mathrm{pH}=-\log[\ce{H3O+}], \qquad
\Delta G^\circ=-RT\ln K, \qquad
E=E^\circ-\frac{0.0592}{n}\log Q
\]
A factor \(10\) in concentration changes pH or Nernst terms by one logarithmic unit. Check whether \(\ln\) or \(\log\) belongs with the numerical constant.

\subsection{Recurring mixed recipes}

\subsubsection{Ideal gas stoichiometry}
\[
n_{\mathrm{gas}}=\frac{pV}{RT}
\]
Convert gas data to moles before applying reaction coefficients. This links gas-law questions with limiting-reactant calculations.

\subsubsection{Percent change in hydronium concentration from pH}
\[
\frac{[\ce{H3O+}]_2}{[\ce{H3O+}]_1}=10^{\mathrm{pH}_1-\mathrm{pH}_2}, \qquad
\%\text{ change}=(\text{ratio}-1)100\%
\]
Use for ocean-acidification or pH-change wording. Lower pH means larger hydronium concentration.

\subsubsection{Recognition keywords, English and Danish}
\[
\begin{array}{c c}
\toprule
\text{exam wording} & \text{go to}\\
\midrule
\text{bundfald, precipitate, sparingly soluble} & K_{sp}, Q_{sp}\\
\text{puffer, buffer, Henderson} & \text{acid/base buffer}\\
\text{galvanisk, electrode, concentration cell} & \text{Nernst/electrochemistry}\\
\text{planar, bond angle, pyramidal} & \text{Lewis/VSEPR}\\
\text{limiting, excess, yield, udbytte} & \text{stoichiometry}\\
\text{vapor pressure, boiling point} & \text{intermolecular forces}\\
\text{rate law, half-life, Arrhenius} & \text{kinetics}\\
\bottomrule
\end{array}
\]
Use the wording to jump to the relevant recipe before doing algebra.
